

\chapter{Convergence}

\marginpar{\red{Lecture 1}}


\section{Types of convergence}

\begin{definition}[Types of Convergence]
    \label{def: types of convergence}
    We say that a sequence of random variables $(X_n)_{n=1}^\infty$ converges to
    a random variable $X$,
    \begin{enumerate}
        \item \textbf{almost surely}, or \textbf{with probability one}, denoted by \(X_n \xrightarrow{a.s.} X\), if
        \[
            \Prob[\Big]{\lim_{n \to \infty} X_n = X} = 1.
        \]

        \item \textbf{in \(L^p\)} for \(p \geq 1\), denoted by \(X_n \xrightarrow{L^p} X\), if
        \(X_n, X \in L^p\) for all \(n\) and
        \[
            \lim_{n \to \infty} \Exp[\Big]{|X_n - X|^p} = 0.
        \]

        \item \textbf{in probability}, denoted by \(X_n \xrightarrow{p} X\), if for all \(\epsilon > 0\)
        \[
            \lim_{n \to \infty} \Prob[\big]{|X_n - X| \ge \epsilon} = 0.
        \]

        \item \textbf{in distribution}, denoted by \(X_n \xrightarrow{d} X\), if
        for all \(x \in \real\) satisfying \(\Prob{X = x} = 0\),
        \[
            \lim_{n \to \infty} \Prob{X_n \leq x} = \Prob{X \leq x}.
        \]
    \end{enumerate}
\end{definition}
\begin{lemma}
    \(X_n \xrightarrow{d} X\) is equivalent to
    \[
        \lim_{n \to \infty} \Prob{X_n < x} = \Prob{X < x} =\Prob{X \le x}.
    \]
    for all \(x\) such that \(\Prob{X=x}=0\).
\end{lemma}
\begin{proof}
    If \(X_n \xrightarrow{d} X\), then we have for any \(x\) with \(\Pr(X=x)=0\)
    \[
        \limsup_{n\to \infty} \Pr(X_n < x) \le \lim_{n\to\infty}\Pr(X_n \le x) = \Pr(X<x)
    \]
    The more difficult part is the lower bound. For this pick a sequence \(\epsilon_k >0\) such that
    \(\epsilon_k \to 0\) and \(\Pr(X = x-\epsilon_k) = 0\). This is possible as
    there is only a countable number of \(y\in \real\) such that \(\Pr(X=y) \neq
    0\) and the complement of a countable set is dense in \(\real\). Then we have
    by \(X_n \xrightarrow{d} X\)
    \[
        \lim_{n\to\infty}\Pr(X_n \le x-\epsilon_k) = \Pr(X \le x - \epsilon_k).
    \]
\end{proof}

\begin{figure}
    \centering
    \begin{tikzcd}[
        % 1. Global Node Styling
        cells={nodes={
            draw, 
            rectangle, 
            rounded corners, 
            align=center, 
            minimum width=2.5cm, 
            inner sep=6pt
        }},
        % 2. Spacing between nodes
        row sep=0.6cm,
        column sep=0.9cm,
        % 3. Global Arrow Styling
    ]
        % --- Row 1 ---
        \parbox{2.1cm}{Almost sure convergence}
        \arrow[r, Rightarrow]{}{\text{\ref{it: as to p}}}
        \arrow[d,xshift=-20]{}{\text{MCT/DCT}}
        & \parbox{2.3cm}{Convergence in probability}
        \arrow[r, Rightarrow] {}{\ref{it: p to d}}
        \arrow[l, out=190, in=350, looseness=1] {}[below,yshift=-2]{\text{Borel-Cantelli}}
        & \parbox{2.4cm}{Convergence in distribution} 
        \arrow[l, out=220, in=340, looseness=1] {}[below,xshift=10,yshift=-2]{X \equiv c \in \real \text{ \ref{it: d to p}}}
        \\
        
        % --- Row 2 ---
        % The loop is defined here on the Lp node
        $L^p$-convergence, $p \ge 1$
        \arrow[r, Rightarrow] {}{\text{\ref{it: Lq to Lp}}}
        \arrow[loop below, Rightarrow, looseness=4, out=220, in=290,xshift=-10]{}[xshift=15,yshift=-3]{L_q \overset{\ref{it: Lq to Lp}}\implies L_p \; (p \le q)}
        & $L^1$ \text{ convergence} \arrow[u, Rightarrow] {} {\text{\ref{it: L1 to p}}}
    \end{tikzcd}
    \caption{
        Implications among different modes of convergence (Thm.~\ref{thm:
        convergence implications}). The Monotone Convergence Theorem (MCT) and
        Dominated Convergence Theorem (DCT) are important results in Measure
        Theory for moving limits into integrals.
    }
\end{figure}

\begin{theorem}[Convergence implications]
    \label{thm: convergence implications}
    The following implications hold:
    \begin{enumerate}[label=(\alph*)]
        \item\label{it: as to p}
        If \(X_n \xrightarrow{a.s.} X\), then \(X_n \xrightarrow{p} X\).
        \item\label{it: Lq to Lp}
        If \(X_n \xrightarrow{L^q} X\) for some \(q \ge p \ge 1\), then \(X_n \xrightarrow{L^p} X\).
        \item\label{it: L1 to p}
        If \(X_n \xrightarrow{L^1} X\), then \(X_n \xrightarrow{p} X\).
        \item\label{it: p to d}
        If \(X_n \xrightarrow{p} X\), then \(X_n \xrightarrow{d} X\).
        \item \label{it: d to p}
        If \(X\) is deterministic, that is \(X \equiv c\in \real\) almost surely, then
        \(X_n \xrightarrow{d} X\) implies \(X_n \xrightarrow{p} X\).
    \end{enumerate}
\end{theorem}
\begin{proof}
    \begin{enumerate}[wide]
        \item[\ref{it: as to p}] The set of point-wise convergence is given by
        \begin{align*}
            \Set{X_n \to X}
            &= \Set[\Big]{\omega \in \Omega : \lim_{n \to \infty} X_n(\omega) = X(\omega)}
            \\
            \overset{\text{def.}}&= \Set[\Big]{
                \omega \in \Omega
                : \forall \epsilon > 0, \exists N\in \nat, \forall n \ge N
                : |X_n(\omega) - X(\omega)| < \epsilon}.
            \\
            &= \bigcap_{\epsilon > 0} \bigcup_{N \in \nat} \bigcap_{n \ge N}
            \Set[\Big]{\omega \in \Omega : |X_n(\omega) - X(\omega)| < \epsilon}.
        \end{align*}
        By de Morgan's laws for set complements and almost sure
        convergence we therefore have
        \begin{align*}
            0 &= \Prob[\big]{X_n \not\to X}
            \\
            &= \Prob[\bigg]{\bigcup_{\epsilon > 0} \bigcap_{N \in \nat} \bigcup_{n \ge N}
            \Set[\Big]{\omega \in \Omega : |X_n(\omega) - X(\omega)| \ge \epsilon}}
            \\
            \overset{\text{fix } \epsilon}&\ge \Prob[\bigg]{\bigcap_{N \in \nat} \bigcup_{n \ge N}
            \Set[\Big]{\omega \in \Omega : |X_n(\omega) - X(\omega)| \ge \epsilon}}
            \\
            &= \lim_{N \to \infty} \Prob[\Big]{\bigcup_{n \ge N}
            \Set[\Big]{\omega \in \Omega : |X_n(\omega) - X(\omega)| \ge \epsilon}}
            \\
            \overset{\text{fix } n=N}&\ge \lim_{N\to \infty} \Prob[\big]{|X_N - X| \ge \epsilon},
        \end{align*}
        where the penultimate equality follows from \(\lim_{N\to \infty} \Prob{A_N} = \Prob[\big]{\bigcap_{N\in \nat} A_N}\).
        for a decreasing sequence of events \(A_1 \supseteq A_2 \supseteq \ldots\).

        \item[\ref{it: Lq to Lp}] See Remark \ref{rem: Lp convergence}.

        \item[\ref{it: p to d}] Let \(x\in \real\) be such that \(\Prob{X=x}=0\).
        If \(X_n \le x\), then either \(X\) is further from \(X_n\) than
        \(\epsilon>0\) or it is also smaller than \(x+\epsilon\).  More
        specifically,
        \begin{equation}
            \label{eq: upper bound for p to d}
            \Set{X_n \le x}
            \subseteq \Set{X \le x + \epsilon}
            \cup \Set[\big]{|X_n - X| > \epsilon}.
        \end{equation}
        With similar arguments we have
        \begin{equation}
            \label{eq: lower bound for p to d}
            \Set{X \le x - \epsilon}
            \subseteq \Set{X_n \le x}
            \cup \Set[\big]{|X_n - X| > \epsilon}
        \end{equation}
        Consequently, we have
        \begin{align*}
            \Prob{X \le x - \epsilon} - \smash{\underbrace{\Prob[\big]{|X_n - X| > \epsilon}}_{\to 0}}
            \overset{\eqref{eq: lower bound for p to d}}&\le \Prob{X_n \le x}
            \\
            \overset{\eqref{eq: upper bound for p to d}}&\le
            \Prob{X \le x + \epsilon} + \underbrace{\Prob[\big]{|X_n - X| > \epsilon}}_{\to 0}.
        \end{align*}
        This implies for all \(\epsilon > 0\)
        \[
            \Prob{X \le x - \epsilon}
            \le \liminf_{n \to \infty} \Prob{X_n \le x}
            \le \limsup_{n \to \infty} \Prob{X_n \le x}
            \le \Prob{X \le x + \epsilon}.
        \]
        Since this is true for all \(\epsilon > 0\) we have
        \[
            \liminf_{n \to \infty} \Prob{X_n \le x}
            \ge \lim_{\epsilon \to 0} \Prob{X \le x - \epsilon}
            = \Prob[\Big]{\bigcap_{\epsilon > 0} \Set{X \le x - \epsilon}}
            = \Prob{X < x},
        \]
        and
        \[
            \limsup_{n \to \infty} \Prob{X_n \le x}
            \le \lim_{\epsilon \to 0} \Prob{X \le x + \epsilon}
            = \Prob[\Big]{\bigcup_{\epsilon > 0} \Set{X \le x + \epsilon}}
            = \Prob{X \le x}.
        \]
        Due to \(\Prob{X \le x} = \Prob{X < x} + \Prob{X = x}\) and \(\Prob{X =
        x} = 0\) we thus have the claim by the squeeze theorem:
        \[
            \lim_{n\to \infty} \Prob{X_n \le x} = \Prob{X \le x}.
        \]
        \item[\ref{it: d to p}] Assume that \(X\equiv c\) and let \(\epsilon>0\).
        Then by the Portemanteau theorem and \(\Pr(X=c-\epsilon) = \Pr(X=c+\epsilon) = 0\)
        we have
        \[
            \Prob[\big]{\abs{X_n -X} < \epsilon}
            = \Pr_{X_n}((c-\epsilon, c+\epsilon))
            \to \Pr_X((c-\epsilon, c+\epsilon)) = 1.
        \]
        and consequently \(\Prob[\big]{\abs{X_n - X} \ge \epsilon} \to 0\).

        \item[\ref{it: L1 to p}] By the \(L^1\) Markov inequality (Corollary
        \ref{cor: markov inequality Lp}), we have for \(\epsilon > 0\) \[
            \Prob[\big]{|X_n - X| > \epsilon}
            \le \frac{\Exp[\big]{|X_n - X|}}{\epsilon}
            \to 0.
            \qedhere
        \]
    \end{enumerate}
\end{proof}

\begin{theorem}[General Markov inequality]
    \label{thm: markov inequality}
    Let \(I\subseteq \real\) be an interval, \(h\colon I\to (0, \infty)\) a
    non-decreasing function and \(X\) an \(I\)-valued random variable.
    Then for all \(t \in I\) we have
    \[
        \Prob{X \ge t} \le \frac{\Exp{h(X)}}{h(t)}.
    \]
\end{theorem}
\begin{corollary}[\(L^p\) Markov inequality]
    \label{cor: markov inequality Lp}
    We have for all \(t > 0\)
    \[
        \Prob[\Big]{|X - \Exp{X}| \ge t}
        \le \frac{\Exp[\Big]{\abs[\big]{X-\Exp{X}}^p}}{t^p}.
    \]
\end{corollary}
\begin{proof}
    Markov inequality applied to \(\tilde X = |X - \Exp{X}|\) 
    and \(h(x) = x^p\) with \(I = [0, \infty)\).
\end{proof}

\begin{proof}[Proof of Theorem \ref{thm: markov inequality}]
    Since \(h\) is non-decreasing and \(t, X \in I\)
    we have that \(X \ge t\) implies \(h(X) \ge h(t)\)
    and thus
    \[\begin{aligned}
        \Prob{X \ge t}
        &\le \Prob[\big]{h(X) \ge h(t)}
        = \Exp{\Indicator{h(X) \ge h(t)}}
        \le \Exp[\Big]{\frac{h(X)}{h(t)}\Indicator{h(X) \ge h(t)}}
        \\
        &\le \Exp[\Big]{\frac{h(X)}{h(t)}},
    \end{aligned}
    \]
    where we have used \(h > 0\) in the last inequality to remove the indicator.
\end{proof}

\subsection{Exercises}

\begin{exercise}[Convergence in distribution but not probability]
    Let \(X\sim \normal{0,1}\) and define
    \[
        X_n \coloneq (-1)^n X
    \]
    Prove that
    \begin{enumerate}[label=(\alph*)]
        \item \(X_n \xrightarrow{d} X\)
        \begin{solution}
            Observe that the distribution of \(X\) is symmetric, that is
            \[\begin{aligned}
                \Pr(-X \le x) &= \Pr(X \ge -x)
                \\
                &= \int_{-x}^\infty \frac{1}{\sqrt{2\pi}} e^{-t^2/2} \, dt
                \\
                \overset{y=-t}&= \int_{-\infty}^x \frac{1}{\sqrt{2\pi}} e^{-y^2/2} \, dy
                \\
                &= \Pr(X \le x).
            \end{aligned}
            \]
            Consequently, for any \(x \in \real\) we have \(\Pr(X_n \le x) = \Pr(X \le x)\).
            In particular this implies \(\Pr(X_n \le x) \to \Pr(X \le x)\) for all \(x\).
        \end{solution}

        \item \(X_n\) does not converge in probability to \(X\).
        \begin{solution}
            We have for \(\epsilon > 0\) and every odd \(n\)
            \[\begin{aligned}
                \Pr(|X_n - X| \ge \epsilon)
                &= \Pr(2\abs{X} \ge \epsilon)
                \\
                &= \Prob[\big]{\abs{X} \ge \epsilon/2}
                \\
                &= 2\int_0^{\epsilon/2} \frac{1}{\sqrt{2\pi}} e^{-t^2/2} \, dt.
                \\
                &\ge 2 \frac{\epsilon}2 \frac{1}{\sqrt{2\pi}} e^{-(\frac{\epsilon}2)^2/2} >0.
            \end{aligned}
            \]
            The sequence \(\Pr(|X_n - X| \ge \epsilon)\) thus does not converge to \(0\).
        \end{solution}
    \end{enumerate}
\end{exercise}

\begin{exercise}[Convergence in probability but not \(L^p\)]
    Let \(U\sim \uniform([0,1))\) be a standard uniform random variable and define
    \[
        X_n \coloneq n \Indicator{[0, 1/n)}(U)
    \]
    Show that
    \begin{enumerate}[label=(\alph*)]
        \item the sequence $X_n$ converges to $0$ in probability,
        \begin{solution}
            We have for any \(\epsilon > 0\)
            \[
                \Pr(|X_n - 0| \ge \epsilon)
                \le \Pr(X_n \neq 0)
                = \Prob[\big]{U\le \tfrac1n}
                = \tfrac1n
                \to 0.
                \qedhere
            \]
        \end{solution}

        \item the sequence \(X_n\) does not converge to \(0\) in \(L^1\),
        \begin{solution}
            We have
            \[
                \Exp{|X_n - 0|}
                = \Exp{X_n}
                = n \cdot \Prob[\big]{U \in [0, 1/n)}
                = n \cdot \frac1n
                = 1 \not\to 0.
                \qedhere
            \]            
        \end{solution}

        \item For every \(p \ge 1\) construct a sequence such that
        \(Y_n\xrightarrow{L^q} X\) for all \(q < p\) but not for \(q \ge p\).
        \begin{solution}
            Define \(Y_n \coloneq X_n^{1/p}\) with \(X_n\) as above. Then
            \[
                \E[|Y_n - 0|^q] = \E[X_n^{q/p}] = n^{q/p} \cdot \Prob[\big]{U \in [0, 1/n)} = n^{q/p - 1}.
            \]
            This converges to \(0\) for all \(q < p\) but not for \(q \ge p\).
        \end{solution}
        \end{enumerate}
\end{exercise}

\begin{exercise}[Convergence in \(L^p\) and probability but not almost surely] 
    Let \(U\sim \uniform([0,1))\) be a standard uniform random variable.
    We will scan the interval for \(U\) with increasingly fine measurements. Specifically,
    let
    \begin{alignat*}{7}
        X_0 &\coloneq \Indicator{[0,1)}(U)
        \\
        X_1 &\coloneq \Indicator{[0,\frac12)}(U)
        \quad & X_2 &\coloneq \Indicator{[\frac12, 1)}(U)
        \\
        X_3 &\coloneq \Indicator{[0, \frac14)}(U)
        \quad & X_4 &\coloneq \Indicator{[\frac14, \frac12)}(U)
        \quad & X_5 &\coloneq \Indicator{[\frac12, \frac24)}(U)
        \quad & X_6 &\coloneq \Indicator{[\frac34, 1)}(U)
    \end{alignat*}
    That is
    \[
        X_n \coloneq \Indicator{\left[
            \frac{k}{2^{m}}, \frac{k+1}{2^{m}}
        \right)}(U),
        \quad \text{with}\quad 
        m \coloneq \lfloor\log_2(n)\rfloor, \quad k \coloneq n - 2^m. 
    \]
    Show that
    \begin{enumerate}[label=(\alph*)]
        \item \(X_n \xrightarrow{L^p} 0\),
        \begin{solution}
            We have for any \(p \ge 1\)
            \[
                \Exp{|X_n|^p} = \int_0^1 \Indicator{
                    \left[\frac{k}{2^{m}}, \frac{k+1}{2^{m}} \right)
                }(u) \, du
                = \frac{1}{2^m} = 2^{-\lfloor \log_2(n) \rfloor} \to 0.
            \]
            Consequently, \(X_n \xrightarrow{L^p} 0\).
        \end{solution}

        \item \(X_n \xrightarrow{p} 0\),
        \begin{solution}
            \(X_n \xrightarrow{L^p} 0\) implies \(X_n \xrightarrow{p} 0\) by Theorem
            \ref{thm: convergence implications} \ref{it: L1 to p}.

            Alternatively:
            \[
                \Prob[\bigl]{|X_n - 0| \ge \epsilon}
                = \Prob{X_n = 1}
                = 2^{-\lfloor \log_2(n) \rfloor} \to 0.
            \]
            for any \(\epsilon \in (0,1]\) and \(\Prob[\bigl]{|X_n - 0| \ge
            \epsilon}=0\) for \(\epsilon > 1\).
        \end{solution}

        \item \(X_n\) does not converge almost surely to \(0\).
        \begin{solution}
            Choose any \(\omega \in \Set{U\in [0,1)}\). We will show that there exists an
            \(\epsilon > 0\), specifically any \(\epsilon < 1\) such that
            \(|X_n(\omega) - 0| \ge \epsilon\) for infinitely many \(n\).
            Since we show this for any \(\omega\in \Set{U\in [0,1)}\), this would imply
            \[
                \Prob{X_n \not\to 0} \ge \Prob{U \in [0,1)} = 1.
            \]
            For any \(N\in \nat\) the task is consequently to find \(n\ge N\)
            such that \(X_n(\omega) = 1\). Let \(m \coloneq \lfloor \log_2(N) \rfloor + 1\).
            Then there exists a \(k \in \{0, \ldots, 2^m-1\}\) such that
            \[
                \frac{k}{2^m} \le U(\omega) < \frac{k+1}{2^m}.
            \]
            Consequently, for \(n = 2^m + k\) we have \(n \ge N\) and
            \[
                X_n(\omega) = \Indicator{\left[
                    \frac{k}{2^{m}}, \frac{k+1}{2^{m}}
                \right)}(U(\omega)) = 1.
                \qedhere
            \]
        \end{solution}
    \end{enumerate}
\end{exercise}



% \begin{exercise}[Polya's urn model]
%     In the strong law of large numbers, the random variable at the limit is a
%     degenerate random variable. A more interesting example is the Polya's urn model.
%     Consider an urn containing $R_0$ red balls and $G_0$ green balls initially. A
%     ball is drawn at random. Thus the drawn ball is red with probability
%     $R_0/(R_0+G_0)$, and is green with probability $G_0/(R_0+G_0)$. The drawn ball
%     is replaced and a ball of the same color is added to the urn. Thus if we denote
%     by $R_1$ and $G_1$ the number of red balls and green balls in the urn after
%     this, $\{R_1=R_0+1,G_1=G_0\}$ with probability $R_0/(R_0+G_0)$, and
%     $\{R_1=R_0,G_1=G_0+1\}$ with probability $G_0/(R_0+G_0)$. It can be shown that
%     as this procedure is repeated, the proportion of red balls converge almost
%     surely and the limiting random variable, say $X$, has distribution
%     $\mathrm{Beta}(R_0,G_0)$. The proportion of green balls converge to $1-X$ and
%     $1-X$ has distribution $\mathrm{Beta}(G_0,R_0)$. Thus in this situation the
%     limit is non-degenerate.
% \end{exercise}

% % \begin{lstlisting}
% % Polya <- function(N,M){
% % r = 1
% % g = 1
% % p = ggplot() +
% %   xlim(0, N) +
% %   ylim(0, 1) +
% %   xlab("n") +
% %   ylab(TeX("$\\frac{R_n}{R_n+G_n}$")) +
% %   ggtitle("Plot of Proportion of Red Balls") +
% %   theme(axis.title.y = element_text(angle = 0, vjust = 0.5)) +
% %   theme(plot.title = element_text(hjust = 0.5))
% % for(m in 1:M){
% %   g.seq = numeric(N)
% %   r.seq = numeric(N)
% %   r.seq[1] = r
% %   g.seq[1] = g
% %   for(i in 2:N){
% %     U = runif(1)
% %     if(U<r.seq[i-1]/(r.seq[i-1]+g.seq[i-1])){
% %       r.seq[i] = r.seq[i-1]+1
% %       g.seq[i] = g.seq[i-1]
% %     }
% %     else
% %     {
% %       r.seq[i] = r.seq[i-1]
% %       g.seq[i] = g.seq[i-1] + 1
% %     }
% %   }
% %   r.prop.seq = r.seq/(r.seq+g.seq)
% %   df = data.frame(x=1:N,y=r.prop.seq)
% %   p = p + geom_line(data=df, aes(x=x,y=y), color=m)
% %   }
% %   return(p)
% % }
% % \end{lstlisting}

% \begin{figure}[h]
%     \centering
%     % \includegraphics[width=\linewidth]{Figures/Rplot02.png}
%     \caption{Plot of $R_n/(R_n+G_n)$ for $n=1,2,\ldots,1000$. Here $R_0=1,G_0=1$. $10$ simulations are shown. In all cases the proportions approach a limit, and the limiting random variable
%     is non-degenerate. The limit is clearly different for the different simulations.}
%     \label{fig:enter-label}
% \end{figure}



\section{\texorpdfstring{\(L^p\)}{Lp} spaces}

\begin{definition}
    We define the \(\mathcal{L}^p\) space to be
    \[
        \mathcal{L}^p \coloneq \Set{X\colon \Omega \to \real : \Exp{\abs{X}^p} < \infty}
    \]
    and define the seminorm for \(X\in \mathcal{L}^p\) as
    \[
        \|X\|_p \coloneq \Exp[\big]{\abs{X}^p}^{\frac1p}.
    \]
\end{definition}

We called \(\|\cdot\|_p\) a ``seminorm''. For this we need to prove that
it actually satisfies the seminorm properties. This is the content of the
following theorem.

\begin{theorem}[Seminorm]
    \label{thm: seminorm properties}
    \(\|\cdot\|_p\) satisfies the
    seminorm properties. That is for \(X,Y \in \mathcal{L}^p\)
    \begin{enumerate}[label=(\roman*)]
        \item\label{it: positive}
        \(\|X\|_p \ge 0\),
        \hfill (positive)
        \item\label{it: scaling}
        \(\|\lambda X\|_p = |\lambda| \|X\|_p\)
        for all \(\lambda \in \real\),
        \hfill (scaling)
        \item\label{it: triangle inequality}
        \(\|X + Y\|_p \le \|X\|_p + \|Y\|_p\).
        \hfill (triangle/Minkowski inequality)
    \end{enumerate}
    Moreover we have for any \(p \le q\)
    \begin{equation}
        \label{eq: Lp norm domination}
        \|X\|_p \le \|X\|_q.
    \end{equation}
\end{theorem}

\begin{remark}[\(L^p\) convergence]
    \label{rem: Lp convergence}
    Observe that \(X_n \xrightarrow{L^p} X\) (Definition \ref{def: types of convergence}) if and only if \(\|X_n -X\|_p \to
    0\). And Equation \eqref{eq: Lp norm domination} means for \(1\le p\le q\) that \(\|X_n - X\|_q \to 0\) implies \(\|X_n - X\|_p \to 0\).
\end{remark}

\begin{example}[Not positive definite]
    Observe that the missing property for \(\|\cdot\|_p\) to be a norm 
    is that \(\|X\|_p = 0\) implies \(X=0\). As an example consider
    a uniformly distributed \(U\sim \uniform(0,1)\) with density \(f(x) = \Indicator{[0,1]}\).
    Define
    \[
        X \coloneq \begin{cases}
            U & U \neq 1
            \\
            0 & U = 1.
        \end{cases}
    \]
    Then clearly \(U\neq X\) but for any \(p\ge 1\)
    \[
        \E[(X-U)^p]
        = \int_0^1 (x^p \Indicator{[0,1)}(x) - x^p) dx
        = \int_0^1 \Indicator{\Set{1}} dx = 0,
    \]
    and consequently \(\|X-U\|_p = 0\).
\end{example}

\begin{fact}
    \label{fact: Lp zero norm}
    \(\|X-Y\|_p = 0\) if and only if \(\Prob{X=Y}=1\).
\end{fact}
\begin{proof}[Proofsketch]
    Measure theory: \(f\ge 0\), \(\int f d\mu = 0\) 
    implies \(\mu(f\neq 0) = 0\).
\end{proof}


\begin{definition}[\(L^p\) space]
    We define the normed space \((L^p,\|\cdot\|_p)\) as
    \[
        L^p \coloneq \Set{[X]: X \in \mathcal{L}^p},
        \quad\text{with}\quad
        \|[X]\|_p \coloneq \|X\|_p
    \]
    where \([X]\) is the equivalence class of the equivalence relation
    \[
        X\sim Y :\iff \|X-Y\|_p = 0 \iff \Pr(X=Y) =1.
    \]
\end{definition}
\begin{remark}[Norm well defined]
    The definition of the norm is well-defined as for \(X,Y\in [X]\)
    we have by the triangle inequality
    \[
        \|X\|_p \le \|Y\|_p + \|X - Y\|_p \overset{X,Y\in [X]}= \|Y\|_p.
    \]
\end{remark}

The key result to prove the seminorm properties in Theorem \ref{thm: seminorm properties}
is the following inequality.

\begin{theorem}[Hölder's inequality]
    \label{thm: Hölder inequality}
    For \(p,q >1\) with \(\frac1p + \frac1q = 1\) we have
    \[
        \Exp{\abs{XY}} \le \Exp{\abs{X}^p}^{\frac1p}\Exp{\abs{Y}^q}^{\frac1q} = \|X\|_p \|Y\|_q
    \]
\end{theorem}
\begin{proof}
    Exercise \ref{ex: hölder's inequality}.
\end{proof}

Using this inequality we can now prove Theorem \ref{thm: seminorm properties}

\begin{proof}[Proof of Theorem \ref{thm: seminorm properties}]
    The properties \ref{it: positive} and \ref{it: scaling} follow
    directly from the definition and are left as an exercise. For
    the norm domination \eqref{eq: Lp norm domination}
    we assume without loss of generality for \(p< q\). We use \(r \coloneq
    \frac{q}p>1\) and \(s \coloneq \frac{p}{q-p}>1\) such that \(\frac1{r} +
    \frac1{s} = 1\). Then
    \[
        \|X\|_p^p = \Exp{\abs{X}^p}
        = \Exp{\abs{\abs{X}^p \cdot 1}}
        \overset{\text{Hölder}}\le \Exp{\abs{X}^{pr}}^{\frac1r} \Exp{1^s}^{\frac1s}
        = \Exp{\abs{X}^q}^{\frac{p}q}
        = \|X\|_q^p.
    \]
    Taking the \(p\)-th root implies the claim.
    What is left is only the triangle inequality \ref{it: triangle inequality},
    which is the most difficult part to prove. Due to this difficulty
    it has its own name the ``Minkowski inequality''. We start with a
    simple application of the triangle inequality for real numbers:
    \[\begin{aligned}
        \|X+ Y\|_p^p
        &= \Exp{\abs{X + Y}^p}
        \\
        &\le \Exp[\big]{\abs{X+Y}^{p-1} \abs{X}} + \Exp[\big]{\abs{X+Y}^{p-1} \abs{Y}}.
    \end{aligned}
    \]
    We then define
    \(q \coloneq \frac{p}{p-1}\) such that \(\frac1p + \frac1q = 1\).
    We may therefore apply Hölder's inequality to both terms to get
    for the first term
    \[
        \Exp[\big]{\abs{X+Y}^{p-1} \abs{X}}
        \overset{\text{Hölder}}\le \Exp[\Big]{\abs{X+Y}^{q(p-1)}}^{\frac1q} \Exp{\abs{X}^p}^{\frac1p}
        = \|X+Y\|_p^{p/q} \|X\|_p,
    \]
    and similarly \(\Exp[\big]{\abs{X+Y}^{p-1} \abs{Y}} \le \|X+Y\|_p^{p/q}
    \|Y\|_p\) for the second.
    Thus
    \[
        \|X+ Y\|_p^p
        \le \|X+Y\|_p^{p/q} (\|X\|_p + \|Y\|_p).
    \]
    Since \(p - \frac{p}q = 1\) we can divide by \(\|X+Y\|_p^{p/q}\)
    to obtain
    \[
        \|X+ Y\|_p = \|X+ Y\|_p^{p-\frac{p}q} \le \|X\|_p + \|Y\|_p.
        \qedhere
    \]
\end{proof}

\begin{fact}[\(L^p\) is a Banach space]
    \label{fact: Lp complete}
    For every \(p \ge 1\) the normed space \((L^p, \|\cdot\|_p)\) is \underline{complete},
    that is every Cauchy sequence converges in \(L^p\). 
\end{fact}
\begin{proof}
See e.g. \citep[Thm.~6.25 and 7.3]{klenkeProbabilityTheoryComprehensive2014}.
\end{proof}

\subsection{Exercises}

\begin{exercise}[Proof of Hölder's inequality]
    \label{ex: hölder's inequality}
    Let \(\frac1p + \frac1q = 1\) with \(p,q > 1\).
    \begin{enumerate}[label=(\alph*)]
        \item Prove ``Young's inequality'': for \(a,b \ge 0\) we have
        \[
            ab \le \frac{a^p}p + \frac{b^q}q.
        \]
        \textbf{Hint:} \emph{Option 1}: Let \(x = p \log(a)\) and \(y = q \log(b)\)
        and use convexity of the exponential function.
        \emph{Option 2 (brute force)}: Find the minimum of 
        the function \(f(x) = \frac{x^p}p + \frac{b^q}q - xb\).
        \begin{solution}
            \emph{Option 1:} Using the convexity of the exponential function,
            \[\begin{aligned}
                ab
                = e^{\log(a) + \log(b)}
                = e^{\frac1p (p\log(a)) + \frac1q (q\log(b))}
                &\le \frac1p e^{p\log(a)} + \frac1q e^{q\log(b)}
                \\
                &= \frac{a^p}p + \frac{b^q}q.
            \end{aligned}
            \]
            \emph{Option 2:} On \([0, \infty)\), \(f\) is convex  as \(f''(x) = (p-1)x^{p-2} \ge 0\).
            The minimum is therefore attained at \(f'(x) = x^{p-1} - b = 0\), that is
            at \(x_* = b^{\frac1{p-1}}\). Since \(q= \frac{p}{p-1}\) we have that \(x_*^p = b^q\).
            \[
                f(x_*) = \frac{b^q}p + \frac{b^q}q - b^{\frac{1}{p-1} + 1} = 0.
            \]
            Consequently, for all \(x \ge 0\) we have \(f(x) \ge f(x_*) = 0\).
            In particular for \(x = a\) we have the claim.
        \end{solution}

        \item Prove ``Hölder's inequality'' (Theorem \ref{thm: Hölder inequality})

        \textbf{Hint:} Apply Young's inequality to \(a =
        \frac{\abs{X}}{\|X\|_p}\) and \(b = \frac{\abs{Y}}{\|Y\|_q}\).
        \begin{solution}
            Applying Young's inequality we have
            \[
                \frac{\abs{X}}{\|X\|_p} \cdot \frac{\abs{Y}}{\|Y\|_q}
                \le \frac1p \Bigl(\frac{\abs{X}}{\|X\|_p}\Bigr)^p
                + \frac1q \Bigl(\frac{\abs{Y}}{\|Y\|_q}\Bigr)^q.
            \]
            Taking expectation on both sides we get
            \[
                \frac{\Exp[\big]{\abs{XY}}}{\|X\|_p \|Y\|_q}
                \le 
                    \frac1p \Exp[\Big]{\Bigl(\frac{\abs{X}}{\|X\|_p}\Bigr)^p}
                    + \frac1q \Exp[\Big]{\Bigl(\frac{\abs{Y}}{\|Y\|_q}\Bigr)^q}
                = 
                    \frac1p + \frac1q
                = 1,
            \]
            where we use that \(\|X\|_p^p = \Exp{\abs{X}^p}\) is deterministic to move
            it outside of the expectation. Multiplying both sides by \(\|X\|_p \|Y\|_q\)
            implies the claim.
        \end{solution}
    \end{enumerate}
\end{exercise}

\begin{exercise}[\(L^2\)-Hilbertspace]
    Show that 
    \[
        \langle [X], [Y] \rangle \coloneq \langle X, Y\rangle \coloneq \Exp{XY}
    \]
    defines an inner product on \(L^2\) which induces the norm \(\|\cdot\|_2\).
    \begin{solution}
        Clearly \(\langle X, X\rangle = \|X\|_2^2\) so the inner product induces
        the norm \(\|\cdot\|_2\). We need to show the inner product properties:
        \begin{enumerate}[label=(\roman*)]
            \item\label{it: inner product positive}
            \(\langle X, X \rangle \ge 0\) with equality if and only if \([X]=0 \overset{\text{Fact } \ref{fact: Lp zero norm}}\iff X=0\) a.s.
            \\ \null \hfill (positive definiteness)
            \item\label{it: inner product symmetry}
            \(\langle X, Y \rangle = \langle Y, X \rangle\).
            \hfill (symmetry)
            \item\label{it: inner product linearity}
            \(\langle aX + bY, Z \rangle = a \langle X, Z \rangle + b \langle Y, Z \rangle\)
            for all \(a,b\in \real\).
            \hfill (linearity in first argument)
        \end{enumerate}
        Property \ref{it: inner product positive} follows from \(\langle X, X\rangle = \|X\|_2^2\)
        and this property is established in Theorem \ref{thm: seminorm properties} and
        Fact \ref{fact: Lp zero norm}.
        Property \ref{it: inner product symmetry} follows from the
        commutativity of multiplication. Property \ref{it: inner product linearity}
        follows from the linearity of expectation.
        \(L^2\) is further a complete normed space by Fact \ref{fact: Lp complete}.
        Consequently, \((L^2, \langle \cdot, \cdot \rangle)\) is a Hilbert space.
    \end{solution}
\end{exercise}


\marginpar{\red{Lecture 2}}

\section{Borel-Cantelli}
\label{sec: borel cantelli}

\subsection{Increasing and Decreasing Sequences of Events}

A sequence of events $A_1,A_2,A_3,\ldots$ is called increasing if 
\[
A_1\subset A_2\subset A_3\subset \ldots.
\]
In this case we write $A_n\uparrow \cup_m A_m$. 

A sequence of events $A_1,A_2,A_3,\ldots$ is called decreasing if 
\[
A_1\supset A_2\supset A_3\supset \ldots.
\]
In this case we write $A_n\downarrow \cap_m A_m$. 

\subsection{Infinitely Often}

Let $A_1,A_2,A_3,\ldots$ be a sequence of events. Consider the event that infinitely many of the events $A_n$ happen. We can represent this event as
\[
\cap_{i=1}^\infty\cup_{j=i}^\infty A_j \;.
\]
Thus, we are saying that for every $u$, there is some $j\geq i$ such that $A_j$ happens. This just means that infinitely many of the events $A_n$ happen. We denote this event by $\limsup A_n$ or $[A_n \mbox{ i.o. }]$.

Similarly we can define the $\liminf$ of a sequence of events.
\[
\liminf A_n = \cup_{i=1}^\infty\cap_{j=i}^\infty A_j \;.
\]
Thus, we are saying that there is some $i$ such that all the events $A_i,A_{i+1},A_{i+2},\ldots$ happens.

So if $\limsup A_n$ does not happen, then $\liminf A_n^c$ happens and vice versa. So $(\limsup A_n)^c = \liminf A_n^c$.


\begin{theorem}[Borel-Cantelli Lemma 1]
Let $\left(A_n\right)_{n=1}^\infty$ be sequence of events. If $\sum_{n=1}^\infty \Prob*{A_n}<\infty$ then $\Prob*{\limsup A_n}=0$.
\end{theorem}


\begin{theorem}[Bore-Cantelli Lemma 2]
Let $\left(A_n\right)_{n=1}^\infty$ be sequence of \textbf{independent} events. If $\sum_{n=1}^\infty \Prob*{A_n}=\infty$, then $\Prob*{\limsup A_n}=1$.
\end{theorem}


\begin{proof}[Borel Cantelli 1]
Let $(A_n)_{n=1}^\infty$ be a sequence of events.\\
Suppose
\[
\sum_{n=1}^\infty \Prob*{A_n} < \infty \;.
\]
Let $A$ denote the event that infinitely many $A_n$s happen. Thus
\[
A = [A_n \mbox{ i.o. }] =
\bigcap_{i=1}^\infty \bigcup_{j=i}^\infty A_i \;.
\]
We need to show that 
\[
\Prob*{A}=0\;.
\]
For $i=1,2,3,\ldots$ let
\[
B_i \coloneqq \cup_{j=i}^\infty A_j\;.
\]
The events $B_i$s are decreasing: $B_{i}$ involves union of more events than $B_{i+1}$. Thus
\[
B_1\supset B_2\supset \cdots \downarrow A\;.
\]
Therefore
\[
\Prob*{A} = \lim_{i\to\infty} \Prob*{B_i}\;.
\]
We have
\[
\Prob*{B_i}\leq\sum_{j=i}^\infty\Prob*{A_j}\;.
\]
Since $\sum_{n=1}^\infty\Prob*{A_n}<\infty$,
we have
\[
\lim_{i\to\infty}\sum_{j=i}^\infty\Prob*{A_j}=0\;.
\]
Therefore,
\[
\lim_{i\to\infty}\Prob*{B_i}=0\;.
\]

\end{proof}


\begin{proof}[Borel Cantelli 2]
Let $(A_n)_{n=1}^\infty$ be a sequence of independent events.\\
This means, for any set of integers $\{n_1,n_2,\cdots,n_k\}$ we have
\[
\Prob*{A_{n_1}\cap A_{n_2}\cap \cdots A_{n_k}}
= \Prob*{A_{n_1}} \cdot \Prob*{A_{n_2}} \cdot \cdots \Prob*{A_{n_k}} \;.
\]
Let $A$ denote the event that infinitely many $A_n$s happen. Thus
\[
A = [A_n \mbox{ i.o. }] = \bigcap_{i=1}^\infty \bigcup_{j=i}^\infty A_i \;.
\]
Suppose 
\[
\sum_{n=1}^\infty \Prob*{A_n} = \infty\;.
\]
We want to show that 
\[
\Prob*{A} = 1\;.
\]
For $i=1,2,3,\ldots$ let
\[
B_i \coloneqq \cup_{j=i}^\infty A_i\;.
\]
Therefore
\[
B_1\supset B_2\supset \cdots \downarrow A\;.
\]
To show $\Prob*{A}=1$ we will show that $\Prob*{A^c}=0$. We observe that
\[
A^c = \bigcup_{i=1}^\infty \bigcap_{j=i}^\infty A_i^c\;.
\]
We also
\[
B_i^c = \cap_{j=i}^\infty A_i^c\;,
\]
and 
\[
B_1^c\subset B_2^c\subset \cdots \uparrow A^c\;.
\]
Therefore
\[
\Prob*{B_i^c}\uparrow\Prob*{A^c}\;.
\]
Therefore, it is enough to show that 
$\Prob*{B_i^c}=0$ for all $i$. We fix an $i$. 
For $m=i,i+1,i+2,\ldots$ let 
\[
C_m \coloneqq \cap_{j=i}^m A_i^c\;.
\]
Thus
\[
C_i \supset C_{i+1} \supset \cdots \downarrow B_i^c\;.
\]
Thus 
\[
\Prob*{B_i^c} = \lim_{m\to\infty} \Prob*{C_m}\;.
\]
By independence of the events $A_i$s we have
\[
\Prob*{C_m} = \prod_{j=i}^m \Prob*{A_i^c} \;.
\]
\begin{tcolorbox}
We are using $A_i^c$s are independent. This follows from the fact that $A_i$s are independent.
\end{tcolorbox}
We observe that
\[
\Prob*{A_i^c} = 1-\Prob*{A_i} \leq e^{-\Prob*{A_i}}\;.
\]
Therefore,
\[
\Prob*{C_m}\leq\exp\left(-\sum_{j=i}^m \Prob*{A_i}\right)\;.
\]
Since $\sum_{n=1}^\infty \Prob*{A_n} = \infty$,
\[
\lim_{m\to\infty}\sum_{j=i}^m\Prob*{A_i}=\infty\;.
\]
Therefore,
\[
\lim_{m\to\infty} \Prob*{C_m} = 0\;.
\]
Therefore, 
\[
\Prob*{B_i^c}=0
\;.
\]
\end{proof}


\subsection{Exercises}

\begin{exercise}[Drunk Men and Drunk Birds]

There is a saying, attributed to  Japanese-American mathematician Shizuo
Kakutani: ``A drunk man will find his way home, but a drunk bird may get lost
forever." Let us try to understand this. A drunk birds can be modeled by three
dimensional random walk. Suppose the bird starts at the origin of the three
dimensional integer lattice $\integer^3$. It moves only at integer time points. For
$n=0,1,2,3,\ldots$ let $S_n=(S_n^x,S_n^y,S_n^z)\in\integer^3$ denote its location.
If at a certain time the bird is at location $(x,y,z)$, then at the next time
point it can be at one of the $8$ locations given by $(x\pm 1,y\pm 1,z\pm 1)$.
Thus we can say

\begin{align*}
S_n^x ={} & S_{n-1}^x + X_n\;,\\
S_n^y ={} & S_{n-1}^y + Y_n\;,\\
S_n^z ={} & S_{n-1}^z + Z_n\;,
\end{align*}
where
\[
X_n = \begin{cases}
1 & \mbox{ with probability } 1/2\;,\\
-1 & \mbox{ with probability } 1/2\;,
\end{cases}
\]
\[
Y_n = \begin{cases}
1 & \mbox{ with probability } 1/2\;,\\
-1 & \mbox{ with probability } 1/2\;,
\end{cases}
\]
\[
Z_n = \begin{cases}
1 & \mbox{ with probability } 1/2\;,\\
-1 & \mbox{ with probability } 1/2\;,
\end{cases}
\]
and $(X_1,X_2,X_3,\ldots,Y_1,Y_2,Y_3,\ldots,Z_1,Z_2,Z_3,\ldots)$ are independent.
\[
\Prob*{S_n=(0,0,0)} 
= \Prob*{S_n^x=0}\Prob*{S_n^y=0}\Prob*{S_n^z=0}
\]
For $n$ odd, 
\[
\Prob*{S_n^x=0}=
\Prob*{S_n^y=0}=
\Prob*{S_n^z=0}=0\;.
\]
For $n$ even, say $n=2N$
\[
\Prob*{S_n^x=0}=
\Prob*{S_n^y=0}=
\Prob*{S_n^z=0}=\left(2^{-2N}\binom{2N}{N}\right)^3\;.
\]
By Stirling's approximation
\[
2^{-2N}\binom{2N}{N} \approx \frac{C}{N^{1/2}} 
\approx \frac{C}{n^{1/2}}.
\]
Thus, for $n$ even
\[
\Prob*{S_n=0} \approx C n^{-3/2}\;.
\]
Thus, the probability that $S_n=0$ happens infinitely often is $0$. Thus, a
drunk bird may get lost forever. But for two dimensional or one dimensional
random walk with probability $1$ $S_n=0$ happens infinitely often. Thus a drunk
man will find his way home. Note that, to prove this we cannot use the second
Borel-Cantelli Lemma. Because the second Borel-Cantelli Lemma requires the
events to be independent, but the events $\{S_n=0\}$ are not independent. But
the result is true and can be proved in different ways, generally you will learn
this in a course on Markov Chains.

\end{exercise}

\section{Law of Large Numbers}


\section{Concentration inequalities}


\section{Continuous mapping}


